% --------------------------------------------------------------------------
\section{Esplorazione LCB per il raffinamento delle stime energetiche}
\label{sec:esplorazione-ucb}
% --------------------------------------------------------------------------
 
Il meccanismo EMA descritto nella sezione precedente consente di
raffinare i profili energetici a runtime, ma presenta un prerequisito
implicito: aggiorna solo le varianti che vengono effettivamente
eseguite. Se la stima iniziale di una variante è sovrastimata rispetto
al suo consumo reale, lo scheduler tenderà a non selezionarla, il
worker non raccoglierà mai campioni su di lei non permettendo la correzione dell'errore. Il problema si 
manifesta in particolare
nelle fasi iniziali del sistema, quando i profili offline — determinati
su un hardware di riferimento — potrebbero non riflettere il
comportamento reale sul nodo edge in produzione, sia per la
\emph{dipendenza dall'hardware} (CPU, workload e condizioni termiche
differenti), sia per la \emph{non stazionarietà} del consumo energetico
(che varia con il carico concorrente, la temperatura del processore e
lo stato della cache del runtime).In queste condizioni, una variante con consumo energetico sovrastimato
tende ad essere sistematicamente penalizzata nella scalarizzazione
Pareto e a non venire mai selezionata — impedendo al meccanismo EMA
di raccogliere campioni reali su di lei e correggere l'errore. È
necessario quindi introdurre un meccanismo che promuova l'esplorazione
di varianti poco osservate, indipendentemente dalla loro stima corrente.
 
\subsection{Il principio di ottimismo sotto incertezza}
\label{subsec:ucb-principio}
 
Il meccanismo adottato si basa sul principio di \emph{ottimismo sotto
incertezza} (Optimism in the Face of Uncertainty), mutuato dalla teoria
dei banditi multi-braccio~\cite{auer2002finite}.Tale meccanismo incentiva
l'esplorazione di varianti sotto-campionate, raccogliendo dati reali e
convergendo progressivamente a stime accurate.
 
Per applicare il principio alla minimizzazione del consumo energetico si
utilizza il Lower Confidence Bound (LCB), equivalente all'UCB classico:
\begin{equation}
    \mathrm{LCB}(i) = \mu_i^\text{EMA}
        - \beta \cdot \frac{\sigma_i}{\sqrt{n_i^\text{eff}}}
    \label{eq:lcb}
\end{equation}
dove $n_i^\text{eff} = n_i + 0.5$, con $n_i$ il numero di campioni
InfluxDB disponibili per la variante $i$. L'addendo costante $0.5$
evita la divisione per zero nel caso $n_i = 0$ e assicura la massima
correzione ottimistica per varianti mai osservate. La
Tabella~\ref{tab:lcb-simboli} descrive i simboli della formula.
 
\begin{table}[t]
    \centering
    \small
    \begin{tabular}{lp{9.5cm}}
        \toprule
        \textbf{Simbolo} & \textbf{Significato} \\
        \midrule
        $\mu_i^\text{EMA}$ &
            Stima EMA dell'energia della variante $i$, letta da etcd
            ($\alpha = 0.2$) \\
        $\sigma_i$ &
            Deviazione standard calcolata sui campioni InfluxDB.\\
        $n_i^\text{eff}$ &
            Numero effettivo di campioni: $n_i + 0.5$ \\
        $\beta$ &
            Parametro di esplorazione (configurabile, default 1.0) \\
        \bottomrule
    \end{tabular}
    \caption{Simboli della formula LCB e relativi valori.}
    \label{tab:lcb-simboli}
\end{table}
 
Una LCB bassa segnala che la variante è sotto-esplorata: il sistema la favorisce
nella selezione per raccogliere dati aggiuntivi. All'aumentare di
$n_i$, la correzione si azzera e la stima converge alla media empirica.
 
 
\subsection{Proprietà di convergenza}
\label{subsec:ucb-convergenza}
 
La formula~\eqref{eq:lcb} gode delle seguenti proprietà asintotiche,
che ne garantiscono il corretto bilanciamento tra esplorazione e
sfruttamento:
 
\begin{itemize}
    \item quando $n_i$ è grande (variante ben esplorata), il termine
    $\sigma_i / \sqrt{n_i^\text{eff}}$ tende a zero e
    $\mathrm{LCB}(i) \to \mu_i^\text{EMA}$: il sistema sfrutta la
    stima corrente senza alterarla;
 
    \item quando $n_i = 0$ (variante mai osservata), la correzione è
    massima ($n_i^\text{eff} = 0.5$)
 
    \item $\beta = 0$ disabilita completamente l'esplorazione,
    riducendo la formula alla sola stima EMA;
 
    \item in caso di indisponibilità di InfluxDB, il sistema degrada a
    $\mu^\text{EMA}$ senza correzione, garantendo il funzionamento
    anche in assenza del componente esterno.
\end{itemize}
 
%Capitolo 5
%In un contesto in cui la stima iniziale $\mu_0$ è artificialmente
%gonfiata rispetto al valore reale $E^*$, la convergenza dell'EMA è
%descrivibile analiticamente: con $\alpha = 0.2$, il valore della stima
%dopo $t$ osservazioni è dato da:
%\begin{equation}
%    \mu^\text{EMA}(t) = E^* + (\mu_0 - E^*) \cdot (1 - \alpha)^t
%    \label{eq:ema-convergenza}
%\end{equation}
%Per $\alpha = 0.2$, la stima raggiunge il 95\% della convergenza in
%circa 15 invocazioni e si stabilizza sul valore reale in 20--25
%invocazioni.
 
